
\documentclass{article}
\usepackage{colortbl}
\usepackage{makecell}
\usepackage{multirow}
\usepackage{supertabular}

\begin{document}

\newcounter{utterance}

\centering \large Interaction Transcript for game `wordle', experiment `high\_frequency\_words\_no\_clue\_no\_critic', episode 5 with reflexion{-}Qwen3.5{-}9B.
\vspace{24pt}

{ \footnotesize  \setcounter{utterance}{1}
\setlength{\tabcolsep}{0pt}
\begin{supertabular}{c@{$\;$}|p{.15\linewidth}@{}p{.15\linewidth}p{.15\linewidth}p{.15\linewidth}p{.15\linewidth}p{.15\linewidth}}
    \# & \multicolumn{2}{c}{Player} && \multicolumn{2}{c}{Game Master} \\
    \hline

    \theutterance \stepcounter{utterance}  
    & & & \multicolumn{4}{p{0.6\linewidth}}{
        \cellcolor[rgb]{0.9,0.9,0.9}{
            \makecell[{{p{\linewidth}}}]{
                \texttt{\tiny{[P1$\langle$GM]}}
                \texttt{You have attempted this same task before and did not succeed. These are the plans you wrote after those attempts.} \\
\texttt{Trial \#1: New plan: I failed because I got lost in a long chain of internal reasoning about letter permutations and failed to output a valid five{-}letter word in the required format. Specifically, I stopped mid{-}thought while analyzing "grit" vs "grits" and never finalized a guess. Next time, I will strictly adhere to the format "explanation: {[}concise logic{]} \textbackslash{}n guess: {[}valid 5{-}letter word{]}". I will ensure the word is exactly five letters long, contains only confirmed letters (g, r, i, t) plus one unknown letter, and respects the positional constraints (t at 5, g not 1, r not 2, i not 3). For instance, I will try "trig" + a vowel or "rig" + a vowel + t, ensuring the final string is a real English word like "trig" (no, 4 letters), "trigs" (no), "grits" (s was red). Let's try "grit" + 'a' {-}> "grita" (not English). "trig" + 's' {-}> "trigs" (no). Wait, letters are g, r, i, t. Need one more. Maybe "trig" is not the set. Re{-}evaluating: g(y), r(y), i(y), s(r), t(g). Word ends in t. Contains g, r, i. Not s. Try "grit" + 'o' {-}> "grito" (no). Try "trig" + 'a' {-}> "triga" (no). Try "rig" + 't' + 'a' {-}> "rigat"? No. Try "trig" + 'e' {-}> "trige" (no). Try "trig" + 'u' {-}> "trigu" (no). Try "trig" + 'y' {-}> "trigy" (no). Try "trig" + 'l' {-}> "trigl" (no). Try "trig" + 'm' {-}> "trigm" (no). Try "trig" + 'p' {-}> "trigp" (no). Try "trig" + 'v' {-}> "trigv" (no). Try "trig" + 'w' {-}> "trigw" (no). Try "trig" + 'x' {-}> "trigx" (no). Try "trig" + 'z' {-}> "trigz" (no). Maybe the unknown letter is 'h'? "grith"? No. 'd'? "grith"? No. 'f'? "grif"? No. 'j'? No. 'k'? No. 'l'? "grilt"? No. 'm'? "grimt"? No. 'o'? "grito"? No. 'p'?} \\
\\ 
\\ 
\texttt{You are a language wizard who likes to guess words by using the given rules.} \\
\\ 
\texttt{Welcome to Wordle! You have six attempts to guess the target word, a valid English word of five lowercase letters (a{-}z). Please use the tags "explanation:" and "guess:" to provide a concise explanation for each guess.} \\
\\ 
\texttt{For instance, if your guess is "apple", your response should be} \\
\texttt{explanation: this is a common five{-}letter English word, and I am starting my guess with this word.} \\
\texttt{guess: apple} \\
\\ 
\texttt{After each guess, your answer will be validated, and you will receive feedback indicating which letters are correct (green), which letters are correct but in the wrong position (yellow), and which letters are incorrect (red). This feedback can be useful in determining which letters to include or exclude in your next guess.} \\
\\ 
\texttt{For example, the feedback for "apple" might be:} \\
\texttt{guess\_feedback: a<yellow> p<yellow> p<green> l<yellow> e<red>} \\
\\ 
\texttt{The explanation should contain details about how the guess\_feedback is used to arrive at a new guess.} \\
\\ 
\texttt{Do not generate any other text.} \\
\\ 
\texttt{Let's begin with your first guess.} \\
            }
        }
    }
     \\ \\

    \theutterance \stepcounter{utterance}  
    & \multicolumn{4}{p{0.6\linewidth}}{
        \cellcolor[rgb]{0.561,0.737,0.561}{
            \makecell[{{p{\linewidth}}}]{
                \texttt{\tiny{[P1$\rangle$GM]}}
                \texttt{explanation: Starting with a common five{-}letter word containing high{-}frequency letters to gather initial feedback.} \\
\texttt{guess: crane} \\
            }
        }
    }
    & & \\ \\

    \theutterance \stepcounter{utterance}  
    & & & \multicolumn{2}{p{0.3\linewidth}}{
        \cellcolor[rgb]{0.9,0.9,0.9}{
            \makecell[{{p{\linewidth}}}]{
                \texttt{\tiny{[GM$|$GM]}}
                \texttt{attempts = 1} \\
\texttt{target\_word = right} \\
\texttt{guess = crane} \\
\texttt{guess\_feedback = c<red> r<yellow> a<red> n<red> e<red>} \\
            }
        }
    }
    & & \\ \\

    \theutterance \stepcounter{utterance}  
    & & & \multicolumn{4}{p{0.6\linewidth}}{
        \cellcolor[rgb]{0.9,0.9,0.9}{
            \makecell[{{p{\linewidth}}}]{
                \texttt{\tiny{[P1$\langle$GM]}}
                \texttt{guess\_feedback: c<red> r<yellow> a<red> n<red> e<red>} \\
\\ 
\texttt{Provide your response only in this format.} \\
\texttt{explanation: details} \\
\texttt{guess: word} \\
            }
        }
    }
     \\ \\

    \theutterance \stepcounter{utterance}  
    & \multicolumn{4}{p{0.6\linewidth}}{
        \cellcolor[rgb]{0.561,0.737,0.561}{
            \makecell[{{p{\linewidth}}}]{
                \texttt{\tiny{[P1$\rangle$GM]}}
                \texttt{explanation: The feedback indicates 'r' is in the word but not in the second position, while 'c', 'a', 'n', 'e' are incorrect. I will try to place 'r' in a different position and include other common vowels and consonants to narrow down the word.} \\
\texttt{guess: right} \\
            }
        }
    }
    & & \\ \\

    \theutterance \stepcounter{utterance}  
    & & & \multicolumn{2}{p{0.3\linewidth}}{
        \cellcolor[rgb]{0.9,0.9,0.9}{
            \makecell[{{p{\linewidth}}}]{
                \texttt{\tiny{[GM$|$GM]}}
                \texttt{attempts = 2} \\
\texttt{target\_word = right} \\
\texttt{guess = right} \\
\texttt{guess\_feedback = r<green> i<green> g<green> h<green> t<green>} \\
            }
        }
    }
    & & \\ \\

    \theutterance \stepcounter{utterance}  
    & & & \multicolumn{2}{p{0.3\linewidth}}{
        \cellcolor[rgb]{0.9,0.9,0.9}{
            \makecell[{{p{\linewidth}}}]{
                \texttt{\tiny{[GM$|$GM]}}
                \texttt{game\_result = WIN} \\
            }
        }
    }
    & & \\ \\

\end{supertabular}
}

\end{document}
